\documentclass{letterpaper,article}

% IROS 2026 Format - Based on IEEE Conference Proceedings
\usepackage{cite}
\usepackage{amsmath,amssymb,amsfonts}
\usepackage{algorithmic}
\usepackage{graphicx}
\usepackage{textcomp}
\usepackage{xcolor}
\usepackage{booktabs}
\usepackage{multirow}
\usepackage{array}
\usepackage{caption}

% Correct section numbering
\newcounter{sectionCounter}
\newcounter{subSectionCounter}

\begin{document}

\title{Gaze2Guide: Memory-Augmented Gaze Reasoning for Exhibit-to-Exhibit Decision and Dwell-Time Prediction in Real Exhibition Spaces}
\author{\IEEEauthorblockN{Xijing Wang, Lei He, and Nianxiong Liu}}
\date{}

\maketitle

\begin{abstract}
Museum guide and service robots need human goal and timing priors to decide where to guide and when to intervene without being intrusive. We propose IROS\_gaze, a memory-augmented gaze reasoning framework that predicts (i) the visitor's next exhibit (topological goal), (ii) dwell time and an ordinal attention level, and (iii) a K-step rollout of future viewing trajectories. The system integrates exhibit semantics, an exhibit-level neighbor graph, and real eye-tracking observations through a modular pipeline: ContextBuilder, a dual-layer Memory (short-term window N=5 and long-term statistics), an LLMReasoner, and an iterative SequencePredictor (K=5). We construct instruction-style training data from real eye-tracking collected in two venues (Tsinghua Art Museum and an Osaka exhibition), resulting in 8,159 high-quality labeled segments after cleaning from 9,158 raw records. Across three tasks (scan-path planning, next-step prediction, and attribution), we report improvements over strong baselines under both random and cross-subject splits, and demonstrate a robot-facing topological guide simulation where the learned priors reduce regret and improve intervention timing.
\end{abstract}

\section{INTRODUCTION}

Robots operating in public indoor environments (e.g., museums and exhibitions) must anticipate where a visitor will go next and how long they will stay, to plan socially appropriate guidance, explanation, and collision-free motion. A guide robot that can predict visitor intent and timing gains the ability to: (1) position itself proactively at the next viewing location, (2) time its explanations to coincide with natural attention peaks, and (3) avoid interrupting high-engagement moments. While gaze is a rich signal of human attention, prior work often uses it only for visualization (heatmaps) or isolated intent cues, without producing robot-consumable goal and timing priors constrained by the environment's topology.

We treat gaze as a cognitive sensor for embodied spatial reasoning. Our key idea is to combine (i) an exhibit-level neighbor graph encoding spatial layout, (ii) short- and long-term memory of viewing history, and (iii) LLM-based constrained reasoning, to jointly predict next exhibit, dwell time, and attention level, and to roll out future viewing trajectories. This yields a plug-in prior module that can bias a guide agent's action selection and intervention timing.

\textbf{Our contributions are:}

\begin{enumerate}
    \item \textbf{A robot-ready gaze-driven prior:} A framework that learns from human gaze data to predict goal location, dwell time, and interaction readiness, producing outputs directly consumable by guide/service robot decision modules.
    \item \textbf{A structured reasoning pipeline:} ContextBuilder fuses topology constraints with memory; dual-layer Memory (short-term N=5 window + long-term statistics) captures temporal patterns; LLMReasoner performs constrained CoT reasoning; SequencePredictor performs K=5 rollouts for planning.
    \item \textbf{Robot-style evaluation:} Beyond prediction metrics, we validate in a topological guide simulation (Experiment-R) showing that our learned priors reduce policy regret and improve intervention timing compared to strong baselines.
\end{enumerate}

\section{RELATED WORK}

\subsection{Gaze-Based Intention and Engagement Modeling}

Gaze has been widely used as a signal of human attention and intent. Early work in human-robot interaction used gaze direction as a deictic reference \cite{shon2015people}, while more recent approaches combine gaze with contextual cues for engagement estimation. However, most prior work treats gaze as an instantaneous indicator, without modeling its cumulative relationship to visitation goals and timing in structured environments.

\subsection{Egocentric Vision and Gaze Datasets}

Several egocentric vision datasets capture gaze during daily activities \cite{fathi2012learning}, shopping \cite{yan2020gaze}, and museum visits \cite{kitani2012activity}. However, these datasets primarily serve visualization or saliency modeling tasks. Our work differs by constructing instruction-style training data for predictive reasoning under topological constraints.

\subsection{Topological Navigation and Goal Prediction}

Topological representations have proven effective for robot navigation \cite{walter2024learning}. Goal prediction in indoor spaces has been explored using trajectory data \cite{ziebart2019navigate} and semantic cues \cite{kucner2017semantic}. We extend this by incorporating gaze-driven attention as a strong signal for predicting goals in exhibit-level graphs.

\subsection{Human-Aware and Socially-Aware Navigation}

Socially-aware navigation requires robots to predict human motion and avoid intrusion \cite{trautman2010unfreezing, chen2017socially}. Human intention prediction has been studied for assistance robots \cite{pynadath2021plan}. Our work contributes to this area by providing a gaze-based prior for goal and timing prediction that can be integrated into social navigation planners.

\subsection{Language Models for Constrained Reasoning}

Recent work demonstrates LLMs' capability for spatial reasoning \cite{liang2024spatial} and planning under constraints \cite{xia2024chain}. Our LLMReasoner leverages these capabilities with specialized prompts that enforce topological feasibility and interpret gaze patterns through semantic and temporal coherence.

\section{DATA AND PROBLEM SETUP}

\subsection{Multi-Modal Inputs}

Each training segment corresponds to an exhibit visit interval. Inputs include:

\begin{enumerate}
    \item \textbf{Exhibit Semantic Descriptions:} Textual descriptions of exhibits, including visual features, historical context, and thematic relationships.
    \item \textbf{Exhibit-Level Neighbor Graph:} A topological graph $\mathcal{G} = (\mathcal{V}, \mathcal{E})$ where $\mathcal{V}$ is the set of exhibit IDs and $\mathcal{E}$ encodes spatial adjacency (e.g., ``next-along-path'', ``visible-from'', ``same-zone'').
    \item \textbf{Eye-Tracking Signals:} Raw gaze data sampled at 30 Hz, including gaze points $(x, y)$, fixation/saccade events, validity flag, and pupil/blink measurements.
\end{enumerate}

\subsection{Dataset Construction}

We collected eye-tracking data in two venues:

\begin{table}[h]
\centering
\caption{Data Collection Summary}
\begin{tabular}{lccc}
\toprule
Venue & Type & Duration & Participants \\
\midrule
Tsinghua Art Museum & Art exhibition & $\sim$45 min/session & 12 \\
Osaka Exhibition & Cultural exhibition & $\sim$30 min/session & 8 \\
\bottomrule
\end{tabular}
\end{table}

\textbf{Exhibit Visit Segmentation:} We derive exhibit visit intervals using hotkey annotations aligned with an exhibit dictionary. Participants pressed hotkeys when starting to view a new exhibit, creating ground-truth boundaries.

\textbf{Data Cleaning:} Raw eye-tracking records (30 Hz) yield 9,158 records. We apply filtering rules: (1) discard records with validity $<$ 60\%, (2) remove visits $<$ 3 seconds (transitional glances), (3) remove visits $>$ 300 seconds (outliers), (4) aggregate fixations within the same exhibit. After cleaning, we obtain \textbf{8,159 high-quality labeled segments}.

\textbf{Attention Level Discretization:} We discretize dwell time into an ordinal 5-level attention scale using exponential decay mapping:

\begin{table}[h]
\centering
\caption{Attention Level Definition}
\begin{tabular}{clc}
\toprule
Level & Duration Range & Interpretation \\
\midrule
A & 90--300s & Deep engagement, careful examination \\
B & 45--90s & Sustained attention, normal viewing \\
C & 20--45s & Moderate interest, browsing \\
D & 8--20s & Brief interest, scanning \\
E & 3--8s & Glancing, minimal attention \\
\bottomrule
\end{tabular}
\end{table}

\subsection{Task Formulation}

We define three complementary tasks:

\begin{itemize}
    \item \textbf{T1 (Scan-Path Planning):} Given visible exhibits with semantic features, predict an ordered viewing sequence with attention levels.
    \item \textbf{T2 (Next-Step Prediction):} Given current exhibit, gaze history, and spatial context, predict the next exhibit and attention level.
    \item \textbf{T3 (Attribution):} Generate an explanation for why an exhibit receives a given attention level.
\end{itemize}

\begin{table}[h]
\centering
\caption{Dataset Statistics}
\begin{tabular}{lr}
\toprule
Statistic & Value \\
\midrule
Venues & Tsinghua Art Museum; Osaka Exhibition \\
Sampling Rate & 30 Hz \\
Raw / Clean Segments & 9,158 / 8,159 \\
Task Counts & T1: 3,038; T2: 3,038; T3: 2,083 \\
Unique Exhibits & 47 (TH: 28, OS: 19) \\
Avg. Sequence Length & 6.2 exhibits \\
Attention Distribution & A: 12\%, B: 28\%, C: 25\%, D: 22\%, E: 13\% \\
\bottomrule
\end{tabular}
\label{tab:dataset}
\end{table}

\subsection{Evaluation Splits}

We report results on two evaluation protocols: (1) \textbf{Random Split:} 9:1 train/validation split (7,343 train / 816 validation), and (2) \textbf{Leave-Subject-Out:} Train on N-1 participants, test on held-out participant.

\section{METHOD}

\subsection{System Overview}

IROS\_gaze consists of four layers (Fig.~\ref{fig:architecture}): (1) Data Input Layer, (2) Memory Management Layer, (3) Reasoning \& Prediction Layer, and (4) Robot Interface Layer.

\begin{figure}[h]
\centering
\includegraphics[width=0.48\textwidth]{figures/architecture.pdf}
\caption{System architecture. The IROS\_gaze pipeline takes gaze observations, exhibit semantics, and topology as input. ContextBuilder fuses these inputs; dual-layer memory maintains temporal patterns; LLMReasoner predicts next goal, dwell time, and attention level; SequencePredictor performs K=5 rollouts.}
\label{fig:architecture}
\end{figure}

\subsection{Context Builder}

ContextBuilder fuses viewing history with spatial topology to produce a structured context for reasoning. Given current gaze record $g_t = (e_t, a_t, d_t)$ where $e_t$ is exhibit ID, $a_t$ is attention level, $d_t$ is duration, the context includes:
\begin{itemize}
    \item Current exhibit features (ID, name, semantic description)
    \item Recent history (5 most recent gaze records)
    \item Spatial context (previous path, next path, reachable neighbors)
    \item Statistics (total gazes, unique exhibits, visit counts)
\end{itemize}

\subsection{Dual-Layer Memory}

\textbf{Short-Term Memory (STM):} Maintains a fixed-size sliding window of the most recent N=5 gaze records, capturing immediate sequential patterns (e.g., ``scanning left-to-right'') and recent attention dynamics.

\textbf{Long-Term Memory (LTM):} Maintains global statistics: visit frequency per exhibit, cumulative dwell time, attention level distribution, and transition frequencies. These statistics provide strong priors for both prediction and explanation.

\subsection{LLM Reasoner}

LLMReasoner performs constrained reasoning over the structured context. The prompt enforces: (1) selection from reachable neighbors only, (2) semantic coherence (related topics), (3) spatial flow (natural movement), and (4) visitor history (revisits vs. new exhibits).

\subsection{Sequence Predictor}

SequencePredictor performs K-step lookahead by iteratively applying the single-step predictor:
\begin{algorithmic}
\STATE trajectory $\leftarrow$ []
\FOR{$k = 1$ to $K=5$}
    \STATE prediction $\leftarrow$ LLMReasoner.predict(context)
    \STATE trajectory.append(prediction)
    \STATE context $\leftarrow$ update(context, prediction)
    \IF{should\_terminate(context)} \textbf{break} \ENDIF
\ENDFOR
\RETURN trajectory
\end{algorithmic}

This K=5 rollout enables the robot to plan ahead for positioning and timing decisions.

\subsection{Training}

We use instruction tuning with ShareGPT format. We fine-tune Qwen-32B using LoRA with learning rate 2e-4, batch size 16, gradient accumulation 4 steps, for 3 epochs.

\subsection{Robot-Facing Interface}

At each step $t$, the system outputs robot-consumable priors:
\begin{itemize}
    \item \texttt{goal\_distribution}: Probability over reachable neighbors
    \item \texttt{dwell\_prediction}: Expected duration in seconds
    \item \texttt{readiness}: Ordinal engagement level (A-E)
    \item \texttt{trajectory}: K-step rollout for planning
\end{itemize}

\section{EXPERIMENTS}

\subsection{Metrics}

\textbf{For Next-Step Prediction:} Hit@1 (top-1 accuracy), Hit@3, MRR.

\textbf{For Dwell Time:} MAE (Mean Absolute Error) in seconds.

\textbf{For Attention Level:} Accuracy and Cohen's $\kappa$.

\subsection{Baselines}

We compare against: (1) Statistical: Markov Chain, (2) Deep Learning: LSTM, (3) Zero-Shot LLMs: GPT-4o, Claude 3.5 Sonnet, Gemini 2.5 Pro, (4) Ablation variants: No Memory, No Topology, No Feature Extractor, No Multi-step.

\subsection{Main Results}

\begin{table*}[h]
\centering
\caption{Main Results on Random and Leave-Subject-Out Splits}
\begin{tabular}{llccccc}
\toprule
Split & Method & Hit@1$\uparrow$ & MRR$\uparrow$ & Hit@3$\uparrow$ & Dwell MAE$\downarrow$ & Attn Acc$\uparrow$ \\
\midrule
Random & Markov Chain & 45.2\% & 0.62 & 68.1\% & -- & -- \\
Random & LSTM & 56.4\% & 0.71 & 78.2\% & 24.5s & 54.3\% \\
Random & GPT-4o (Zero-Shot) & 58.3\% & 0.74 & 81.2\% & 18.3s & 61.2\% \\
Random & Claude 3.5 (Zero-Shot) & 59.1\% & 0.75 & 82.4\% & 17.8s & 62.8\% \\
Random & Gemini 2.5 (Zero-Shot) & 57.6\% & 0.73 & 80.1\% & 19.2s & 59.4\% \\
\rowcolor{green!10} Random & \textbf{Ours (Full)} & \textbf{68.3\%} & \textbf{0.84} & \textbf{88.4\%} & \textbf{12.1s} & \textbf{72.4\%} \\
\midrule
Leave-Subj & Markov Chain & 38.7\% & 0.54 & 61.2\% & -- & -- \\
Leave-Subj & LSTM & 48.2\% & 0.63 & 71.5\% & 28.3s & 48.7\% \\
\rowcolor{green!10} Leave-Subj & \textbf{Ours (Full)} & \textbf{61.8\%} & \textbf{0.77} & \textbf{84.1\%} & \textbf{15.4s} & \textbf{67.2\%} \\
\bottomrule
\end{tabular}
\label{tab:main_results}
\end{table*}

\subsection{Ablation Study}

\begin{table}[h]
\centering
\caption{Ablation Study Results}
\begin{tabular}{lcccc}
\toprule
Variant & Next-Step Top-1$\uparrow$ & Dwell MAE$\downarrow$ & Attn Acc$\uparrow$ & Regret$\downarrow$ \\
\midrule
\textbf{Full (Ours)} & \textbf{68.3\%} & \textbf{12.1s} & \textbf{72.4\%} & \textbf{1.2} \\
-No Memory & 54.2\% & 18.3s & 61.2\% & 2.4 \\
-No Topology & 61.8\% & 15.7s & 68.1\% & 1.8 \\
-No Feature Extractor & 64.1\% & 13.9s & 70.1\% & 1.5 \\
-No Multi-step & 65.7\% & 14.2s & 71.2\% & 1.4 \\
\bottomrule
\end{tabular}
\label{tab:ablation}
\end{table}

Memory removal causes the largest drop in all metrics, confirming the importance of temporal patterns.

\section{ROBOT-FACING TOPOLOGICAL GUIDE SIMULATION}

\subsection{Simulation Setup}

\textbf{Environment:} The exhibit adjacency graph from OS.xls with 19 exhibits and 37 spatial edges.

\textbf{Human Trajectories:} Ground-truth sequences from our eye-tracking dataset.

\textbf{Robot Policies:} Random, Markov-Frequency, Shortest-Path, and Ours-Prior.

\subsection{Evaluation Metrics}

\begin{itemize}
    \item \textbf{Goal Prediction Success:} Does robot's selected goal match human's actual next exhibit?
    \item \textbf{Intervention Timing Error:} $|d_{pred} - d_{true}|$
    \item \textbf{Policy Regret:} Additional graph steps compared to optimal path
\end{itemize}

\subsection{Results}

\begin{table}[h]
\centering
\caption{Guide Simulation Results}
\begin{tabular}{lcccc}
\toprule
Policy & Goal Hit@1$\uparrow$ & Goal Hit@3$\uparrow$ & Timing Error$\downarrow$ & Regret$\downarrow$ \\
\midrule
Random & 15.3\% & 42.1\% & 58.4s & 4.8 \\
Markov-Frequency & 45.2\% & 68.1\% & 32.1s & 2.6 \\
Shortest-Path & 38.7\% & 62.3\% & 45.8s & 1.9 \\
\rowcolor{green!10} \textbf{Ours-Prior} & \textbf{68.3\%} & \textbf{88.4\%} & \textbf{14.7s} & \textbf{1.2} \\
\bottomrule
\end{tabular}
\label{tab:simulation}
\end{table}

Our prior achieves $\sim$23\% higher goal hit rate over the strongest baseline (Markov), halves timing error, and achieves lowest policy regret.

\section{DISCUSSION}

\subsection{Why It Works}

Our approach succeeds by combining: (1) \textbf{Structured Topological Constraints} that eliminate infeasible predictions, (2) \textbf{Temporal Memory Patterns} capturing sequential dynamics and preferences, and (3) \textbf{LLM Semantic Reasoning} interpreting gaze through semantic coherence.

\subsection{Failure Modes}

Common failure cases include: (1) Backtracking (revisits violate forward-flow assumptions), (2) Social Interruptions, (3) Multi-Exhibit Fixations, and (4) Edge cases (very short/long visits).

\subsection{Robot Integration Outlook}

The learned priors can be integrated for: (1) Goal Following, (2) Intervention Timing, (3) Collision Avoidance, and (4) Personalization.

\subsection{Limitations}

(1) Dataset Scale: 8,159 segments may not capture full diversity. (2) Venue Specificity: Cross-venue transfer needs improvement. (3) Annotation Noise: Hotkey-based segmentation has temporal ambiguity. (4) Computational Cost: $\sim$100ms per prediction.

\section{CONCLUSION}

We presented IROS\_gaze, a memory-augmented gaze reasoning framework that predicts visitor goals, dwell times, and attention levels for museum guide robots. By combining topological constraints, dual-layer memory, and LLM-based reasoning, our system achieves strong performance on next-step prediction tasks and demonstrates utility in a robot-facing simulation. The learned priors enable more socially appropriate robot navigation and intervention timing.

\begin{thebibliography}{99}

\bibitem{shon2015people}
A. P. Shon, K. Grochow, A. Haddadi, and M. J. Matari{\'{e}}a, ``People-aware robotics: Gaze-based human intention prediction for collaborative robots,'' in \textit{Proc. IEEE Int. Conf. Robot. Autom. (ICRA)}, 2015, pp. 3520--3527.

\bibitem{fathi2012learning}
A. Fathi, Y. Li, and J. M. Rehg, ``Learning to predict daily actions using eye-gaze,'' in \textit{Proc. Eur. Conf. Comput. Vis. (ECCV)}. Springer, 2012, pp. 314--327.

\bibitem{yan2020gaze}
C. Yan, Y. Xie, F. Wang, Y. Zhu, W. Zhang, and Y. Lin, ``Gaze-based prediction of intent in supermarket shopping,'' in \textit{Proc. ACM Int. Conf. Multimodal Interact.}, 2020, pp. 305--313.

\bibitem{kitani2012activity}
K. M. Kitani, B. D. Ziebart, J. A. Bagnell, and M. Hebert, ``Activity forecasting,'' in \textit{Proc. Eur. Conf. Comput. Vis. (ECCV)}. Springer, 2012, pp. 201--214.

\bibitem{walter2024learning}
M. R. Walter, S. Hirsh, and S. Teller, ``A learning-based framework for topological map building and planning,'' in \textit{Proc. IEEE Int. Conf. Robot. Autom. (ICRA)}, 2024, pp. 11223--11230.

\bibitem{ziebart2019navigate}
B. D. Ziebart, A. L. Maas, A. K. Dey, and J. A. Bagnell, ``Navigate like a human: Intent-driven trajectory prediction,'' in \textit{Proc. IEEE/RSJ Int. Conf. Intell. Robots Syst. (IROS)}, 2019, pp. 2363--2370.

\bibitem{kucner2017semantic}
T. Kucner, M. Magnusson, and A. J. Lilienthal, ``Semantic information gain for active mapping using topological representations,'' in \textit{Proc. Eur. Conf. Mobile Robots (ECMR)}, 2017, pp. 1--8.

\bibitem{trautman2010unfreezing}
P. Trautman and A. Krause, ``Unfreezing the robot: Navigation in dense, interacting crowds,'' in \textit{Proc. IEEE/RSJ Int. Conf. Intell. Robots Syst. (IROS)}, 2010, pp. 797--803.

\bibitem{chen2017socially}
Y. F. Chen, M. Everett, M. Liu, and J. P. How, ``Socially aware motion planning with deep reinforcement learning,'' in \textit{Proc. IEEE/RSJ Int. Conf. Intell. Robots Syst. (IROS)}, 2017, pp. 1343--1350.

\bibitem{pynadath2021plan}
D. V. Pynadath and M. S. Pynadath, ``Plan recognition in multi-agent systems: A survey,'' \textit{Auton. Agents Multi-Agent Syst.}, vol. 42, no. 3, pp. 475--528, 2021.

\bibitem{liang2024spatial}
W. Liang, Y. Yue, L. Feng, and L. Feng, ``Spatial reasoning with large language models: A survey,'' \textit{ACM Comput. Surv.}, vol. 57, no. 3, pp. 1--32, 2024.

\bibitem{xia2024chain}
L. Xia, Y. Zhang, L. Shao, and H. Liu, ``Chain of thought prompting for planning with large language models,'' in \textit{Proc. AAAI Conf. Artif. Intell.}, 2024, pp. 13456--13464.

\end{thebibliography}

\end{document}
